\documentclass[12pt]{article}
\usepackage[utf8]{inputenc}
\usepackage[T1]{fontenc}
\usepackage[margin=1in]{geometry}
\usepackage{setspace}
\usepackage{xcolor}
\usepackage[colorlinks=true, urlcolor=blue!60!black, linkcolor=blue!60!black]{hyperref}
\setlength{\parindent}{2em}
\setlength{\parskip}{1.25em}
\renewcommand{\baselinestretch}{1.0}
\title{\textit{Central Hudson Gas \& Electric Corp. v. Public Service Commission}, 447 U.S. 557 (1980)}
\author{Case Brief}
\date{}
\begin{document}
\maketitle

\section{Case Metadata}
\begin{description}
  \item[Case Name \& Citation:] \textit{Central Hudson Gas \& Electric Corp. v. Public Service Commission of New York}, 447 U.S. 557 (1980) (No.~79-565).
  \item[Argued / Decided:] Argued March 17, 1980; Decided June 20, 1980.
  \item[Vote:] 8--1 (reversing the New York Court of Appeals, \textit{Consolidated Edison Co. v. Public Service Comm'n}, 47 N.Y.2d 94, 390 N.E.2d 749 (1979)).
  \item[Author:] Justice Lewis F. Powell, Jr. (joined by Chief Justice Burger and Justices Stewart, White, and Marshall).
  \item[Concurrences / Dissent:] Justice Brennan (concurring in the judgment, 447 U.S. at 572, joining Justices Blackmun and Stevens); Justice Blackmun (concurring in the judgment, joined by Justice Brennan, 447 U.S. at 573); Justice Stevens (concurring in the judgment, joined by Justice Brennan, 447 U.S. at 579); Justice Rehnquist (dissenting, 447 U.S. at 583).
  \item[Chapter / Topic:] Free Speech \& Press --- Commercial Speech and Intermediate Scrutiny (U.S. Const. amends. I, XIV).
  \item[Primary Source URLs:] \href{https://supreme.justia.com/cases/federal/us/447/557/}{Justia U.S. Supreme Court Center}; \href{https://www.law.cornell.edu/supremecourt/text/447/557}{Cornell LII}.
\end{description}

\section{Facts and Procedural History}
In December 1973, during the 1973--1974 Mideastern oil embargo and fuel shortage, the New York Public Service Commission (PSC) ordered electric utilities in New York to cease all advertising ``promot[ing] the use of electricity,'' citing insufficient fuel stocks to meet winter customer demand. Three years later, after the immediate shortage subsided, the PSC solicited public comments on retaining the ban. Central Hudson Gas \& Electric Corp., a regulated electric utility holding a geographic monopoly in its service area, opposed continuation on First Amendment grounds.

On February 25, 1977, the PSC issued a Policy Statement extending the prohibition indefinitely. It divided utility advertising into two categories: (1) \emph{promotional advertising} (``intended to stimulate the purchase of utility services''), which it banned completely as contrary to national energy conservation policy; and (2) \emph{institutional and informational advertising} (not clearly intended to promote sales, such as encouraging consumers to shift usage from peak to off-peak hours), which it permitted subject to PSC review. In denying rehearing, the PSC added that promotional advertising would worsen rate inequities because New York electricity rates were not based on marginal cost, causing all ratepayers to subsidize new demand, and would send ``misleading signals'' during an energy crisis.

Central Hudson sued in New York state court under the First and Fourteenth Amendments. The trial court, the Appellate Division, and the New York Court of Appeals upheld the order, reasoning that promotional advertising by a regulated monopoly in a ``noncompetitive market'' carried minimal informational value and was outweighed by the State's interest in energy conservation. The Supreme Court noted probable jurisdiction and reversed.

\section{Issue(s) Presented}
\begin{enumerate}
  \item Does a state utility commission's complete ban on promotional advertising by a regulated electric utility restrict commercial speech protected by the First and Fourteenth Amendments?
  \item What standard of judicial scrutiny governs restrictions on truthful, nonmisleading commercial speech, and does New York's total ban satisfy that standard?
\end{enumerate}

\section{Holding}
Yes on protection, and no on constitutionality (8--1). Although the Constitution accords commercial speech lesser protection than other guaranteed expression, restrictions on truthful commercial speech concerning lawful activity must satisfy a four-part intermediate scrutiny test. New York's complete ban fails the fourth prong because it is more extensive than necessary to serve the State's substantial interest in energy conservation.

\section{Reasoning}
Writing for a five-Justice majority, Justice Powell synthesized the Court's post-\textit{Virginia Pharmacy Board v. Virginia Citizens Consumer Council}, 425 U.S. 748 (1976), jurisprudence into the canonical four-part \textit{Central Hudson} test (447 U.S. at 566):
\begin{enumerate}
  \item \textbf{Protected Commercial Speech:} The expression must ``concern lawful activity and not be misleading.'' Because commercial speech derives its First Amendment value from its informational function---and is hardier and more verifiable than political speech---false, deceptive, or unlawful commercial messages receive no protection.
  \item \textbf{Substantial Governmental Interest:} The asserted governmental interest must be ``substantial.''
  \item \textbf{Direct Advancement:} The regulation must ``directly advance[] the governmental interest asserted'' rather than offering only ineffective or remote support.
  \item \textbf{Narrow Tailoring:} The restriction must be ``not more extensive than is necessary to serve that interest.''
\end{enumerate}

Applying the four prongs to New York's ban:
\begin{itemize}
  \item \textbf{Prong 1 (Satisfied):} Central Hudson's promotional ads were neither misleading nor unlawful. Rejecting the New York Court of Appeals' monopoly rationale, Powell noted that electric utilities still face interfuel competition from oil and natural gas suppliers, and even monopoly advertising gives consumers useful information about new services or usage levels.
  \item \textbf{Prong 2 (Satisfied):} Both asserted interests---energy conservation and fair, efficient utility rate structures---were ``clear and substantial.''
  \item \textbf{Prong 3 (Split):} The rate-equity justification failed because the causal chain between promotional advertising, off-peak usage, peak demand, and rate subsidies was ``tenuous'' and ``highly speculative.'' By contrast, the energy-conservation interest passed because an ``immediate connection'' exists between utility advertising and electricity demand.
  \item \textbf{Prong 4 (Failed):} The total ban was fatally overbroad. It barred Central Hudson from advertising energy-efficient electric devices---such as heat pumps or electric backup for solar heating---that would cause no net increase in total energy use or would displace less efficient fuels. Moreover, the PSC failed to show that narrower measures (such as requiring disclosures about relative efficiency and cost, or previewing ad campaigns) would inadequately advance conservation.
\end{itemize}

\section{Concurring and Dissenting Opinions}
\begin{itemize}
  \item \textbf{Justice Blackmun (Concurring in the Judgment, joined by Justice Brennan, 447 U.S. at 573--579):} Criticized the majority's four-part test as too weak. While intermediate scrutiny fits anti-deception or time-place-manner rules, suppressing truthful information about a lawful product merely to ``dampen'' demand is covert paternalism that ``strikes at the heart of the First Amendment.'' If New York wanted to curb energy use, it had to regulate consumption directly (e.g., thermostat limits) rather than ``keeping the public in ignorance.''
  \item \textbf{Justice Stevens (Concurring in the Judgment, joined by Justice Brennan, 447 U.S. at 579--583):} Faulted Powell's first definition of commercial speech (``expression related solely to the economic interests of the speaker and its audience'') as dangerously overbroad, sweeping in labor strikes or economic commentary. Stevens argued New York's ban suppressed core public-issue advocacy---such as arguing that electric heat is cleaner than wood stoves---and thus violated the strict standard of \textit{Whitney v. California}, 274 U.S. 357 (1927).
  \item \textbf{Justice Brennan (Concurring in the Judgment, 447 U.S. at 572--573):} Joined both Stevens (on the breadth of the ban) and Blackmun (on rejecting manipulation of public conduct by suppressing truthful commercial speech).
  \item \textbf{Justice Rehnquist (Dissenting, 447 U.S. at 583--606):} Mounted a sweeping critique of commercial-speech doctrine on three grounds: (1) a state-created monopoly subject to plenary rate regulation should not claim First Amendment immunity from conservation orders; (2) elevating commercial advertising to near-parity with political speech resurrects \textit{Lochner v. New York}, 198 U.S. 45 (1905), turning ordinary economic regulation into constitutional litigation; and (3) the ``no more extensive than necessary'' prong invites judges to second-guess complex utility engineering (such as whether heat pumps boost summer air-conditioning peaks).
\end{itemize}

\section{Key Quotations}
\begin{quote}
``The Constitution therefore accords a lesser protection to commercial speech than to other constitutionally guaranteed expression.'' \\
--- Justice Powell, 447 U.S. at 562--563.
\end{quote}

\begin{quote}
``In commercial speech cases, then, a four-part analysis has developed. At the outset, we must determine whether the expression is protected by the First Amendment. For commercial speech to come within that provision, it at least must concern lawful activity and not be misleading. Next, we ask whether the asserted governmental interest is substantial. If both inquiries yield positive answers, we must determine whether the regulation directly advances the governmental interest asserted, and whether it is not more extensive than is necessary to serve that interest.'' \\
--- Justice Powell, 447 U.S. at 566.
\end{quote}

\begin{quote}
``No differences between commercial speech and other protected speech justify suppression of commercial speech in order to influence public conduct through manipulation of the availability of information.'' \\
--- Justice Blackmun, concurring in the judgment, 447 U.S. at 578.
\end{quote}

\begin{quote}
``The Court in so doing returns to the bygone era of \textit{Lochner v. New York}, 198 U.S. 45 (1905), in which it was common practice for this Court to strike down economic regulations adopted by a State based on the Court's own notions of the most appropriate means for the State to implement its considered policies.'' \\
--- Justice Rehnquist, dissenting, 447 U.S. at 589.
\end{quote}

\section{Doctrinal Significance and Modern Legacy (Through 2026)}
\textit{Central Hudson} remains the governing doctrinal framework for commercial speech, though its rigor has tightened significantly since 1980 (see \href{https://www.law.cornell.edu/constitution-conan/amendment-1/central-hudson-test-and-current-doctrine}{Cornell LII Constitution Annotated}):
\begin{itemize}
  \item \textbf{Tailoring Refinement and Anti-Paternalism:} In \textit{Board of Trustees of SUNY v. Fox}, 492 U.S. 469 (1989), the Court clarified that Prong 4 does not demand the ``least restrictive means,'' only a ``reasonable fit'' proportionate to the interest. Yet in \textit{44 Liquormart, Inc. v. Rhode Island}, 517 U.S. 484 (1996), the Court invalidated a liquor price-advertising ban, rejecting the ``vice'' exception of \textit{Posadas de Puerto Rico Associates v. Tourism Co.}, 478 U.S. 328 (1986), and echoing Blackmun's concurrence that paternalistic bans on truthful commercial speech demand ``special care.'' Justice Thomas urged abandoning \textit{Central Hudson} altogether in favor of strict scrutiny whenever government suppresses truthful ads to keep consumers ignorant.
  \item \textbf{Heightened Scrutiny for Speaker- and Content-Based Rules:} In \textit{Sorrell v. IMS Health Inc.}, 564 U.S. 552 (2011), the Court struck down Vermont's ban on selling pharmacy prescriber records for pharmaceutical marketing, holding that content- and speaker-based burdens on commercial speech trigger ``heightened judicial scrutiny'' before finding the law failed \textit{Central Hudson} as well.
  \item \textbf{Compelled Disclosures (\textit{Zauderer}) vs. Restrictions (\textit{Central Hudson}) Through 2026:} For \emph{compelled} commercial disclosures of ``purely factual and uncontroversial information,'' \textit{Zauderer v. Office of Disciplinary Counsel}, 471 U.S. 626 (1985), applies relaxed rational-basis review (unless unjustified or unduly burdensome, as in \textit{NIFLA v. Becerra}, 585 U.S. 755 (2018)). Through 2025--2026, the boundary between \textit{Zauderer}, \textit{Central Hudson}, and strict scrutiny has become the central battleground in platform and ESG litigation: in \href{https://uk.practicallaw.thomsonreuters.com/Document/I6dc2e9a0580011f0aa648a444928c7d7/View/FullText.html}{\textit{DoorDash, Inc. v. City of New York} (S.D.N.Y. June 2025)}, a federal court struck down a municipal customer-data sharing mandate under \textit{Central Hudson} for speculative interests and poor fit; and in \href{https://cdn.ca9.uscourts.gov/datastore/opinions/2026/11/09/25-5327.pdf}{\textit{U.S. Chamber of Commerce v. Sanchez} (9th Cir. Nov. 9, 2026)}, the Ninth Circuit analyzed California's corporate climate-disclosure statutes (SB 253 and SB 261) across the \textit{Zauderer}/\textit{Central Hudson} divide, contrasting objective commercial disclosures with subjective, predictive, or politically fraught mandates.
\end{itemize}

\section{Pedagogical Critiques and Exam Hypotheticals}
\textbf{Pedagogical Critiques:}
\begin{itemize}
  \item \textit{The Boundary Problem:} Powell offered two inconsistent definitions of commercial speech---``expression related solely to the economic interests of the speaker and its audience'' (p.~561) versus ``speech proposing a commercial transaction'' (p.~562). Subsequent cases (e.g., \textit{Bolger v. Youngs Drug Products Corp.}, 463 U.S. 60 (1983)) treat the ``proposal of a commercial transaction'' as the core, but mixed political-commercial corporate speech (such as ESG pledges) still blurs the line.
  \item \textit{First Amendment Lochnerism:} Rehnquist's dissent presaged modern critiques that aggressive commercial-speech review (\textit{Sorrell}, \textit{44 Liquormart}) allows courts to deregulate consumer protection, data privacy, and public-health regimes under the guise of free speech.
\end{itemize}

\textbf{Exam Hypothetical:} To curb youth nicotine addiction and carbon emissions from lithium batteries, State X enacts a statute banning all outdoor billboard and social-media advertisements for flavored disposable e-cigarettes, including truthful ads targeted at adult smokers comparing vaping's lower combustion risk to traditional cigarettes. A licensed manufacturer challenges the ban. \textit{Analysis:} Under \textit{Central Hudson}, (1) the ads concern a lawful product for adults and are truthful (protected); (2) reducing youth addiction and e-waste are substantial interests; (3) the State must produce non-speculative evidence---not mere conjecture---that banning ads \emph{directly} reduces youth use or waste; and (4) a total ban reaching adult-oriented truthful comparisons likely fails Prong 4 (fit) under \textit{Lorillard Tobacco Co. v. Reilly}, 533 U.S. 525 (2001), and \textit{44 Liquormart}, because age-gated digital targeting, mandatory health warnings (\textit{Zauderer}), or direct sales/battery-recycling regulations would advance the State's goals without suppressing truthful adult information.

\textbf{``Opposite Side'' Contrast:} Contrast \textit{Central Hudson} with \textit{Zauderer v. Office of Disciplinary Counsel}, 471 U.S. 626 (1985), and \textit{Consolidated Edison Co. v. Public Service Commission}, 447 U.S. 530 (1980) (decided the same day). Whereas \textit{Central Hudson} applies intermediate scrutiny to government \emph{suppression} of commercial advertising, \textit{Zauderer} applies deferential review to compelled \emph{disclosures} of purely factual, uncontroversial commercial terms (reasoning that disclosure \emph{adds} to the stream of consumer information rather than closing it). Conversely, \textit{Consolidated Edison} struck down the same PSC's ban on utility bill inserts addressing nuclear power under full \emph{strict scrutiny} because the utility spoke on a controversial public issue rather than proposing a commercial transaction.

\end{document}
